\documentclass{article}

% --- Encoding & Language ---
\usepackage[utf8]{inputenc}     % pdfLaTeX
\usepackage[T1]{fontenc}

% --- Math & Chemistry ---
\usepackage{amsmath, amsthm, amssymb}

\title{Introductory Real Analysis: Square Roots - Exam Problems}
\author{Robin Rovenszky}
\date{Last updated: April 15, 2026}

\newtheorem{theorem}{Theorem}
\newtheorem{remark}{Remark}
\newtheorem{proposition}{Proposition}
\newtheorem{definition}{Definition}
\newtheorem{example}{Example}

\begin{document}

\maketitle
\tableofcontents
\newpage

\section{Exam Problems}
\subsection{Exam Beta}
\subsubsection{Problem 1.}
Evaluate the following: (3 points)
\vspace{8pt}
    \hrule
    \vspace{8pt}
    \[
    \sqrt{\frac{36}{25}\cdot\frac{81}{144}\cdot0,36}
    \]
    First, let $0,36 = \frac{36}{100}$, then split the square root into a product of fractions
    \[
    =\sqrt{\frac{36}{25}}\cdot \sqrt{\frac{81}{144}}\cdot\sqrt{\frac{36}{100}} = \frac{6}{5}\cdot\frac{9}{12}\frac{6}{10}=\frac{6\cdot9\cdot6}{5\cdot12\cdot10}=\frac{324}{600}=\frac{27}{50}=0,54
    \]
    \vspace{3pt}
    \hrule 
    \vspace{8pt}
    \[
    \sqrt{72 \cdot 32}
    \]
    \[
    \sqrt{2^8 \cdot 3^2} = 2^4 \cdot 3 = 16 \cdot 3 = 48.
    \]
    Generally,
    \[
    \sqrt{ab}
    \]
    Eval. $ab$'s prime factorization
    \[
    ab = p_1^{\alpha_1} \cdot p_2^{\alpha_2} ... p_i^{\alpha_i}
    \]
    If all $\alpha_i$ are even, then
    \[
    \sqrt{ab} = p_1^{\frac{\alpha_1}{2}} \cdot p_2^{\frac{\alpha_2}{2}} ... p_i^{\frac{\alpha_i}{2}}
    \]
    \vspace{3pt}
    \hrule
    \vspace{8pt}
    \[
    \sqrt{90\cdot6,4}
    \]
    \[
    \sqrt{90 \cdot \frac{64}{10}}
    \]
    \[
    \sqrt{10 \cdot 9 \cdot \frac{64}{10}}
    \]
    \[
    \sqrt{9 \cdot 64}
    \]
    \[
    \sqrt{9} \cdot \sqrt{64} = 3 \cdot 8 = 24.
    \]

\subsubsection{Problem 2.}
Evaluate the following expressions (2 points)
\[
2\sqrt{72} - \sqrt{50} - 2\sqrt{8}
\]
\[
\sqrt{72} = \sqrt{36 \cdot 2} = 6\sqrt{2}
\]
\[
\sqrt{50} = \sqrt{25 \cdot 2} = 5\sqrt{2}
\]
\[
\sqrt{8} = \sqrt{4 \cdot 2} = 2\sqrt{2}
\]
\[
2(6\sqrt{2}) - 5\sqrt{2} - 2(2\sqrt{2}) = 12\sqrt{2} - 5\sqrt{2} - 4\sqrt{2} = 3\sqrt{2}
\]

\vspace{3pt}
\hrule
\vspace{8pt}

\[
\sqrt{12} - \sqrt{27} + \frac{1}{2}\sqrt{48}
\]
We notice that all of the numbers are divisible by 3.
\[
2\sqrt{3}  - 3\sqrt{3} + \frac{1}{2}4\sqrt{3}
\]
\[
2\sqrt{3} - 3\sqrt{3} + 2\sqrt{3}
\]
\[
\sqrt{3}.
\]

\subsubsection{Problem 3.}
Simplify the following expressions. (2 points)
\[
\sqrt{30}\cdot\sqrt{20}
\]
\[
\sqrt{30 \cdot 20} = \sqrt{600}
\]
\[
600 = 2^3 \cdot 3 \cdot 5^2
\]
\[
=\sqrt{2^2 \cdot 5^2 \cdot 2 \cdot 3} = 2 \cdot 5 \cdot \sqrt{6} = 10\sqrt{6}
\]

\vspace{3pt}
\hrule
\vspace{8pt}

\[
\sqrt{63}\cdot\sqrt{28}
\]
\[
\sqrt{1764} = ?.
\]
\[
1764 = 2^2\cdot3^2\cdot7^2
\]
\[
\sqrt{2^2\cdot3^2\cdot7^2} = 2\cdot3\cdot7 = 42.
\]

\subsubsection{Problem 4.}
Evaluate the expression. (3 points)
\[
(2\sqrt{18} + \sqrt{54})\cdot(3\sqrt{6}-\sqrt{27}-\sqrt{3})
\]
Simplify every radical
\[
2\sqrt{18} = 2\cdot3\sqrt{2}=6\sqrt{2}
\]
\[
\sqrt{54} = 3\sqrt{6}
\]
\[
\sqrt{27} = 3\sqrt{3}
\]
Then the expression becomes
\[
(6\sqrt{2} + 3\sqrt{6}) \cdot (3\sqrt{6} - 3\sqrt{3} - \sqrt{3})
\]
\[
(6\sqrt{2} + 3\sqrt{6}) \cdot (3\sqrt{6} - 4\sqrt{3})
\]
After simply algebraically expanding the brackets like with any expression, we get
\[
54 + 36\sqrt{3} - 36\sqrt{2} - 24\sqrt{6}.
\]

\subsubsection{Problem 5.}
Calculate the exact value of the following expression. (3 points)
\[
(\sqrt{2 + 2\sqrt{3}} - \sqrt{-2 +2\sqrt{3}})^2
\]
We will omit the details here, ask any LLM.
The final result, for checking purposes, is $4(\sqrt{3} - \sqrt{2})$.

\subsubsection{Problem 6.}
Calculate the exact value of the following expression. (3 points)
\[
\sqrt{48  + 6\sqrt{15}} + \sqrt{48 - 6\sqrt{15}}
\]
We will omit the details here, ask any LLM. The final results, for checking purposes, is $6\sqrt{5}.$

\subsubsection{Problem 7.}
Rationalize the denominator.

\[
\frac{2}{\sqrt{5}-2}
\]
The conjugate of $\sqrt{5}-2$ is $\sqrt{5} + 2$, therefore we will multiply the entire fraction by that.
\[
\frac{2(\sqrt{5} + 2)}{(\sqrt{5}-2)(\sqrt{5} + 2)}
\]
\[
2\sqrt{5} + 4.
\]

\vspace{3pt}
\hrule
\vspace{8pt}

\[
\frac{28}{1+\sqrt{5}}
\]
Multiply by the conjugate
\[
\frac{28(1-\sqrt{5})}{(1+\sqrt{5})(1-\sqrt{5})}
\]
We omit the small details. The final result is
\[
7(\sqrt{5} - 1).
\]

\subsubsection{Problem 8.}
Simplify the following expression. (3 points)
\[
\sqrt{a + 2 + 4\sqrt{a-2}}
\]

\subsection{Exam Alpha}
\subsubsection{Problem 1.}

Calculate the value of the following expressions (? points)

\[
\sqrt{\frac{64}{225} \cdot \frac{100}{9} \cdot 2,25} = \frac{8}{15} \cdot \frac{10}{3} \cdot \frac{15}{10} = \frac{8}{3}
\]

\vspace{3pt}
\hrule
\vspace{8pt}

\[
\sqrt{75 \cdot 48} = \sqrt{3 \cdot 25 \cdot 3 \cdot 16} = 3 \cdot 5 \cdot 4 = 60
\]

\vspace{3pt}
\hrule
\vspace{8pt}

\[
\sqrt{2,5 \cdot 14,4} = \sqrt{\frac{25}{10} \cdot \frac{144}{10}} = \frac{5 \cdot 12}{10} = 6
\]

\subsubsection{Problem 2.}
Evaluate the following expressions. (? points)

\[
\sqrt{80} + \frac{1}{2}\cdot\sqrt{20} + 3\sqrt{45} 
\]

\[
4\sqrt{5} + \frac{1}{2}2\sqrt{5} + 3\cdot 3\cdot\sqrt{5}
\]

\[
4\sqrt{5} + \sqrt{5} + 9\sqrt{5}
\]

\[
14\sqrt{5}
\]

\vspace{3pt}
\hrule
\vspace{8pt}
Evaluate the following expressions

\[
5\sqrt{3} - \frac{1}{3} \cdot \sqrt{27} - \sqrt{48}
\]

\[
5\sqrt{3} + \frac{1}{3} \cdot 3 \cdot \sqrt{3} - 4\sqrt{3}
\]

\[
2\sqrt{3}.
\]

\subsubsection{Problem 3.}
1 point
\[
\sqrt{50} \cdot \sqrt{125} = 5\cdot\sqrt{2} \cdot 5 \sqrt{5} = 25\sqrt{10}
\]

\[
\sqrt{48} \sqrt{27} = 4\sqrt{3} \cdot 3 \cdot \sqrt{3} = 36
\]

\subsubsection{Problem 4.}
\[
(3\sqrt{3} - \sqrt{5})(\sqrt{20}+2\sqrt{12} + 2\sqrt{3}) = (3\sqrt{3} - \sqrt{5})(2\cdot\sqrt{5} + 2\cdot2\cdot\sqrt{3} + 2\sqrt{3}) = (3\sqrt{3} - \sqrt{5})(6\sqrt{3} + 2\sqrt{5}) = 2 \cdot (3\sqrt{3} - \sqrt{5})(3\sqrt{3} + \sqrt{5}) = 2\cdot (27 - 5) = 44.
\]

\subsubsection{Problem 5.}
\[
(\sqrt{6-\sqrt{11}} + \sqrt{6 + \sqrt{11}})^2 = 6 - \sqrt{11} + 6 + \sqrt{11} + 2\sqrt{(6 - \sqrt{11})(6 + \sqrt{11})} = 12 + 2 \sqrt{36 - 11} = 12 + 10 = 22.
\] % to check that this is correct

\subsubsection{Problem 6.}
\[
\sqrt{17 - 4\sqrt{15}} + \sqrt{17 + 4\sqrt{15}}
\]
Then $17$ is $a^2 + b^2$ and $4\sqrt{15}$ is $2ab$, therefore $ab = 2\sqrt{15}$
Now we can make a table
a | 2 | 2\sqrt{3} | 2 \sqrt{5} | 1
b | \sqrt{15} | \sqrt{5} | \sqrt{3} | 2\sqrt{15}
a^2 + b^2 | 19 | 17 | 23 | 61
Now we can continue our solution
\[
2\sqrt{3} + \sqrt{5} + 2\sqrt{3} - \sqrt{5} = 4 \sqrt{3}
\]

\subsubsection{Problem 7.}
\[
\frac{2}{\sqrt{3} - 2} \cdot \frac{\sqrt{3} + 2}{\sqrt{3} + 2} = \frac{2\sqrt{3} + 4}{3 - 4} = -2\sqrt{3} - 4
\]

\[
\frac{12}{\sqrt{17 + 5}} \cdot \frac{5 - \sqrt{17}}{5 - \sqrt{17}} = \frac{12(5-\sqrt{17)}}{25-17} = \frac{}
\]

\subsubsection{Problem 8.}
\[
\sqrt{3 - \sqrt{5}}
\]

We notice that
\[
\frac{6 - 2\sqrt{5}}{2} = \frac{(\sqrt{5} - 1)^2}{2}
\]

Now we can continue our proof
\[
\sqrt{\frac{(\sqrt{5} - 1)^2}{2}} = \frac{\sqrt{5} - 1}{\sqrt{2}} = \frac{\sqrt{10} - \sqrt{2}}{2}
\]

\end{document}